\clearpage\begingroup\SourceLeaf{604b}
\ArithmeticHead{49}{R}{X. CHRONOLOGISCHE ARBEITEN.}
\renewcommand{\thefootnote}{\arabic{footnote})}\setcounter{footnote}{1}
\begin{center}\large X.\par\textls[120]{Chronologische Arbeiten.}\end{center}
\noindent\hspace*{\parindent}In der Augustnummer 1800 der Monatlichen Correspondenz zur Beförderung der Erd- und Himmelskunde, herausgegeben vom Frhn. v. \textsc{Zach},
findet sich eine Veröffentlichung von \textsc{Gauss} über eine Methode zur Berechnung des Osterfestes\footnote{Werke VI, S. 73.}. Die Entdeckung fand wenige Monate vorher statt;
dies zeigt die Notiz Nr. 107 im \textit{Tagebuch}\footnote{Werke X, 1, S. 547.}:

»Iisdem diebus circa (Mai. 16. [1800]) problema chronologicum de festo
paschali eleganter resolvimus.«
\par\noindent X 2 Abh. 6.\hfill 7
\clearpage\SourceLeaf{605}\ArithmeticHead{50}{L}{PHILIPP MAENNCHEN, GAUSS ALS ZAHLENRECHNER.}
\setcounter{footnote}{0}
\noindent\hspace*{\parindent}Wie man aus der Einleitung zu dem genannten Aufsatz ersieht, handelt
es sich bei dem chronologischen Problem nicht um eine astronomische, sondern um eine zahlentheoretische Aufgabe. \textsc{Gauss} schreibt: »Die Absicht dieses
Aufsatzes ist nicht, das gewöhnliche Verfahren zur Bestimmung des Osterfestes zu erörtern, das man in jeder Anweisung zur mathematischen Chronologie findet, und das auch an sich leicht genug ist, wenn man einmal die
Bedeutung und den Gebrauch der dabei üblichen Kunstwörter, \textls[120]{güldene Zahl},
\textls[120]{Epakte}, \textls[120]{Ostergrenze}, \textls[120]{Sonnenzirkel} und \textls[120]{Sonntagsbuchstabe} weiss
und die nötigen Hilfstafeln vor sich hat: sondern von dieser Aufgabe eine
von jenen Hilfsbegriffen unabhängige und bloss auf den einfachsten Rechnungsoperationen beruhende rein analytische Auflösung zu geben.« Die Frage
liegt nahe, welche äussere Veranlassung ihn auf diese Aufgabe gebracht
haben mag.

Fand er das Problem in der Literatur? Die Möglichkeit liegt vor, da bereits
\textsc{Joh. H. Lambert} eine solche analytische Auflösung versucht hatte. Diese Lösung
findet sich im Astronomischen Jahrbuch für das Jahr 1778 unter dem Titel:
\textit{Einige Anmerkungen über die Kirchenrechnung}. \textsc{Lambert} berechnet der Reihe nach
die güldene Zahl, daraus die Epakte, daraus die Ostergrenze, ganz nach den
kirchlichen Vorschriften in voller Umständlichkeit. Daran schliesst sich eine
Rechenvorschrift, um von der Ostergrenze zum nächsten Sonntag zu gelangen.
Auch diese Vorschrift ist zwar, sobald sie erläutert ist, rein analytisch, aber
sie ist nicht praktisch, und vor allem führt von ihr aus kein gangbarer Weg
zur gregorianischen Reform. Das hat \textsc{Lambert} selbst eingesehen; denn er
sagt in der genannten Arbeit im Hinblick auf den gregorianischen Kalender
folgendes:

»Eine allgemeine Formel ist teils an sich zu weitläufig, teils hört sie insofern auf, allgemein zu sein, als \textsc{Clavius}, der den gregorianischen Kalender
eingerichtet hat, von den dabei zu Grunde gelegten allgemeinen Regeln selbst
einige Male abgewichen.«

Es ist nicht ausgeschlossen, dass \textsc{Gauss} diese Arbeit gelesen hat und dadurch angeregt wurde, etwas Besseres zu schaffen. Aber viel wahrscheinlicher ist eine Vermutung, die \textsc{Brendel}\footnote{Mündliche Mitteilung.} aufstellt: \textsc{Gauss} wünschte das genaue
\clearpage\SourceLeaf{606}\ArithmeticHead{51}{R}{X. CHRONOLOGISCHE ARBEITEN.}
\setcounter{footnote}{0}
\noindent Datum seiner Geburt zu ermitteln, über das wohl im Elternhause nichts Urkundliches vorhanden war\footnote{Die auf die Geburt von \textsc{Gauss} bezügliche Eintragung in dem Kirchenbuch von St. Katharinen
zu Braunschweig ist in H. \textsc{Macks} Schrift \textit{C. F. Gauss und die Seinen} (Braunschweig 1927) im Faksimile
wiedergegeben; sie lautet:
\par\noindent\hspace*{\parindent}»Mstr. \textsc{Gebhard Diterich Gass}, Bürger und Gassenschlächter hat mit seiner Ehefrau \textsc{Dorothea} geb. \textsc{Benzen} einen Sohn gezeuget \{den 30 sten April\} dessen Gevattern sind 1) \textsc{Christine
Margaretha Friderica Sieversen} 2) H. \textsc{Johann Gottlieb Wagenknecht} 3) Mons. \textsc{Georg
Carl Ritter}. Das Kind heisst:
\par\centerline{\textsc{Johann Friderich Carl}«}
\par\noindent\hspace*{\parindent}Das hier in \{Klammern\} gesetzte Datum ist in der Handschrift mit kleinerer Schrift über den eigentlichen Text geschrieben; \textsc{Mack} sagt darüber (a. a. O. S. X), der damalige Kirchenbuchführer habe anfänglich den Geburtstag anzugeben vergessen und ihn, wenn auch noch selber, allem Anschein nach erst nach
einiger Zeit nachgetragen.}, und seine Mutter konnte ihm nur angeben, er
sei 1777 am Mittwoch vor Rogate geboren. Es hat viel Wahrscheinlichkeit
für sich, dass \textsc{Gauss} die ihm zugänglichen Anweisungen und Hilfstafeln benutzte, und dass ihm bei dieser Gelegenheit der Gedanke kam, diese Hilfsmittel durch eine auf zahlentheoretischem Wege abzuleitende Formel entbehrlich zu machen.

Auf die Begründung seiner Osterformel hat \textsc{Gauss} verzichtet. Die Unterdrückung des Beweises entschuldigt er mit den Worten\footnote{Werke VI, S. 75.}: »Die Analyse, vermittelst welcher obige Formel gefunden wird, beruhet eigentlich auf Gründen
der höheren Arithmetik, in Rücksicht auf welche ich mich gegenwärtig noch
auf keine Schrift beziehen kann, und lässt sich daher freilich in ihrer ganzen
Einfachheit hier nicht darstellen .\ .\ ..«

Will man den Spuren seiner Gedankengänge nachgehen, so sind zwei
Wege denkbar: 1) Zahlentheoretischer Anschluss an die chronologischen Festsetzungen bei den Bestimmungen des Konzils von Nicäa, bezw. bei der gregorianischen Reform dieser Bestimmungen. 2) Aufstellung einer Formel durch
Beobachtung der vorliegenden Tabellen für die Osterdaten und Ostergrenzen\footnote{Mit dem Ausdruck »Ostergrenze« bezeichnet man in der Chronologie das Datum des Frühlingsvollmonds.}.

Es ist sehr wahrscheinlich, dass \textsc{Gauss} zunächst den zweiten Weg eingeschlagen hat, und dass er nur hie und da nachträglich den ersten wählte,
um seine Folgerungen beweiskräftig zu gestalten. Wir wollen jetzt versuchen,
einen solchen Weg zu rekonstruieren.

Zunächst soll eine Formel für den julianischen Kalender hergeleitet
\par\hfill $7^{\ast}$
\clearpage\SourceLeaf{607}\ArithmeticHead{52}{L}{PHILIPP MAENNCHEN, GAUSS ALS ZAHLENRECHNER.}
\noindent werden, die alsdann für den gregorianischen Kalender passend umgearbeitet
wird. Als Beobachtungsmaterial dienen Tabellen für Osterdaten und Ostergrenzen, die zweifellos \textsc{Gauss} zugänglich waren.

Ein flüchtiger Überblick über eine Ostertabelle liefert leicht die Beobachtung, dass alle Osterdaten des Julianischen Kalenders nach 532 Jahren in
der gleichen Reihenfolge wiederkehren. Daraus ergibt sich, dass das Osterdatum einzig und allein von den Resten der Jahreszahl $(\mathrm{mod}\ 19)$, $(\mathrm{mod}\ 4)$ und
$(\mathrm{mod}\ 7)$ abhängig ist.

Bezeichnet man diese Reste der Reihe nach mit $a$, $b$, $c$, so muss sich
das Osterdatum als eine Funktion von $a$, $b$ und $c$ ergeben.

Ganz besonders einfach ist die Tabelle der Ostergrenzen. Diese Daten
kehren bereits nach 19 Jahren in der gleichen Reihenfolge wieder und befolgen ein unmittelbar erkennbares Gesetz, nämlich:

Beginnen wir den 19jährigen Zyklus mit einem Jahr, das durch 19 geteilt den Rest Null lässt, so ist für dieses Jahr die Ostergrenze am 5. April,
den ich im folgenden als 36. März bezeichnen will. Für jedes folgende Jahr
sind entweder 19 Tage zu addieren, oder 11 Tage zu subtrahieren (da eine
obere und eine untere Grenze für den Frühlingsvollmond festgesetzt ist), um
die Ostergrenze für das betreffende folgende Jahr zu erhalten. Die Tabelle
für einen solchen Zyklus, den man als den \textls[120]{Metonischen Zyklus} bezeichnet,
lautet also:
\begin{center}\small\setlength{\tabcolsep}{6pt}
\begin{tabular}{c|c|c||c|c|c}
Rest $(\mathrm{mod}\ 19)$ & Ostergrenze & $d$ & Rest $(\mathrm{mod}\ 19)$ & Ostergrenze & $d$\\
0&36&15&10&46&25\\
1&25&4&11&35&14\\
2&44&23&12&24&3\\
3&33&12&13&43&22\\
4&22&1&14&32&11\\
5&41&20&15&21&0\\
6&30&9&16&40&19\\
7&49&28&17&29&8\\
8&38&17&18&48&27\\
9&27&6&&&
\end{tabular}\end{center}
\noindent\hspace*{\parindent}Beim Übergang zum neuen Zyklus werden 12 Tage subtrahiert (statt 11),
damit man wieder zum Anfangsdatum zurückkommt. Die Einrichtung beruht
\clearpage\SourceLeaf{608}\ArithmeticHead{53}{R}{X. CHRONOLOGISCHE ARBEITEN.}
\noindent darauf, dass 19 Jahre nahezu ein ganzzahliges Vielfaches des synodischen
Monats sind. Ostern fällt nach den Beschlüssen des Konzils von Nicäa auf
den ersten Sonntag nach dem Frühlingsvollmond. Demgemäss setzt \textsc{Gauss}
das Osterdatum $O$ zusammen aus der Ostergrenze, die er mit $d+21$ bezeichnet,
und der Anzahl der Tage bis zum nächsten Sonntag, $e+1$. Somit erhält er:
\[
O=22+d+e.
\]
\noindent Durch diese Bezeichnungsweise wird erreicht, dass $d<30$ und $e<7$ ist.
Für $d=0$ erhält man den 21. März als Ostergrenze, und diese findet man
in der obigen Tabelle bei $a=15$ als niedrigste. Für $d=29$ würde man
den 50. März erhalten; unsere Tabelle zeigt jedoch, dass der 49. März (für
$a=7$) die äusserste Ostergrenze ist. $e+1$ wurde gewählt, da man für den
Fall, dass die Ostergrenze auf einen Sonntag fällt, 7 Tage addieren muss,
während $e$ nur bis 6 anwachsen kann.

Wir haben in der Tabelle (S. 52) neben die Spalte für die Ostergrenzen
noch eine Spalte für die zugehörigen Werte von $d$ gesetzt. Man erkennt
leicht, dass diese Werte auch unabhängig von denen der Ostergrenze aus dem
ersten Wert $d=15$ abgeleitet werden können, indem man jedesmal 19 addiert und von dem Ergebnis den Rest $(\mathrm{mod}\ 30)$ nimmt. Daher ist
\[
d\equiv19a+k\; (\mathrm{mod}\ 30),
\]
\noindent und aus $d=15$ für $a=0$ erkennt man, dass $k\equiv15$ sein muss. Daher die
Formel:
\[
d\equiv19a+15\; (\mathrm{mod}\ 30).
\]

Wir untersuchen nun die Abhängigkeit der Grösse $e$ von den übrigen
eingeführten Grössen $a$, $b$, $c$, $d$. Von den Grössen $a$ und $d$ können wir eine
ausschalten, da durch die obige Formel die eine mit der anderen verknüpft
ist. Wir schalten $a$ aus; denn wir können die endgültige Osterformel:
\[
O=22+d+e
\]
\noindent dazu benutzen, um die Abhängigkeit von $d$ und $e$ zu ermitteln. Wir gehen
zu diesem Zweck von der Tatsache aus, dass in dem Jahre $J+28$ jedes Datum
wieder auf denselben Wochentag fällt, wie in dem Jahre $J$. Folglich muss
die Kongruenz bestehen:
\[
O_J\equiv O_{J+28}\; (\mathrm{mod}\ 7).
\]
\clearpage\SourceLeaf{609}\ArithmeticHead{54}{L}{PHILIPP MAENNCHEN, GAUSS ALS ZAHLENRECHNER.}
\noindent Da $28=4\mathord{.}7$ ist, so haben $b$ und $c$ für $J$ die gleichen Werte, wie für $J+28$
und natürlich auch für $J+\lambda\mathord{.}28$ $(\lambda=1,2,3,\mathord{.}\mathord{.}\mathord{.})$.

Für die Jahre $J+\lambda\mathord{.}28$ ist also der Siebenerrest von $O_{J+\lambda\mathord{.}28}$ eine von
$a$ und $d$ unabhängige Grösse, sobald man den Siebenerrest von $O_J$, den wir
mit $\nu$ bezeichnen wollen, festgestellt hat. Man hat also:
\[
22+d+e\equiv\nu\; (\mathrm{mod}\ 7),
\]
\noindent woraus sich ergibt:
\[
e\equiv\nu-22-d\qquad(\mathrm{mod}\ 7)
\]
\noindent oder
\[
e\equiv\nu+6+6d\; (\mathrm{mod}\ 7).
\]
\noindent Hierbei ist $e$ nur von den Werten von $b$ und $c$ für das Jahr $J$ abhängig und
sicher von $d$ unabhängig. Es zeigt sich also:
\begin{center}$e$ ist eine lineare Funktion von $d$.\end{center}
\noindent\hspace*{\parindent}Nun wollen wir eine Gruppe von 28 aufeinanderfolgenden Jahren zusammenstellen, einen sogenannten Sonnenzyklus. Und zwar beginnen wir mit
einem Jahre, für das $b=0$ ist. Neben jedes Jahr schreiben wir die Wochentagsnummer irgend eines Tages; die Wochentagsnummern der Tabelle beziehen
sich dann natürlich alle auf den nämlichen Tag.
\begin{center}\footnotesize\setlength{\tabcolsep}{4pt}\renewcommand{\arraystretch}{1.15}
\begin{tabular}{@{}c||c||c||c@{}}
\begin{tabular}[t]{@{}c|c@{}}
&\shortstack{Wochentags-\\Nummer}\\
$J+0$&$n$\\
$J+1$&$n+1$\\
$J+2$&$n+2$\\
$J+3$&$n+3$\\
$J+4$&$n+5$\\
$J+5$&$n+6$\\
$J+6$&$n$\\
$J+7$&$n+1$
\end{tabular}&
\begin{tabular}[t]{@{}c|c@{}}
&\shortstack{Wochentags-\\Nummer}\\
$J+8$&$n+3$\\
$J+9$&$n+4$\\
$J+10$&$n+5$\\
$J+11$&$n+6$\\
$J+12$&$n+1$\\
$J+13$&$n+2$\\
$J+14$&$n+3$\\
$J+15$&$n+4$
\end{tabular}&
\begin{tabular}[t]{@{}c|c@{}}
&\shortstack{Wochentags-\\Nummer}\\
$J+16$&$n+6$\\
$J+17$&$n$\\
$J+18$&$n+1$\\
$J+19$&$n+2$\\
$J+20$&$n+4$\\
$J+21$&$n+5$\\
$J+22$&$n+6$
\end{tabular}&
\begin{tabular}[t]{@{}c|c@{}}
&\shortstack{Wochentags-\\Nummer}\\
$J+23$&$n$\\
$J+24$&$n+2$\\
$J+25$&$n+3$\\
$J+26$&$n+4$\\
$J+27$&$n+5$
\end{tabular}
\end{tabular}\end{center}
\noindent Man erkennt, dass die Wochentagsnummer um 1 wächst, wenn der ganzzahlige
Bestandteil des 4. Teils der neuen Jahreszahl unverändert bleibt, dagegen um
2, wenn dieser Bestandteil um 1 wächst. Die Wochentagsnummer des Jahres
$J+\lambda$ gestaltet sich daher für $\lambda=4k+\nu$ $(\nu=0,1,2,3)$ folgendermassen:
\[
n_{J+\lambda}\equiv n_J+\lambda+k\; (\mathrm{mod}\ 7).
\]
\noindent Wir stellen jetzt die folgende Gruppe auf:
\clearpage\SourceLeaf{610}\ArithmeticHead{55}{R}{X. CHRONOLOGISCHE ARBEITEN.}
\begin{center}
$\left.\begin{array}{l}J\\J+133\\J+266\\J+399\end{array}\right\}$
\begin{minipage}{.39\linewidth}Für diese Jahre haben\newline
$a$ und $c$ die gleichen\newline Werte. Auch $d$ hat\newline den gleichen Wert.\end{minipage}
\end{center}
\noindent\hspace*{\parindent}Da $133\equiv21\; (\mathrm{mod}\ 28)$ ist, so folgt:
\[
n_{J+133}\equiv n_{J+21}\equiv n_J+21+5\equiv n_J+5\; (\mathrm{mod}\ 7).
\]
\noindent Dabei ist allerdings noch vorausgesetzt, dass $b_J=0$ ist.

Da $d$ innerhalb der gewählten Gruppe konstant ist. so ist
\[
d_{J+133}=d_J.
\]
\noindent Nun ist
\[
n_{J+133}\equiv n_J+5,
\]
\noindent daher
\[
e_{J+133}\equiv e_J-5\equiv e_J+2\; (\mathrm{mod}\ 7).
\]
\noindent D. h.: Wächst $b$ von 0 auf 1 (da $133\equiv1\; (\mathrm{mod}\ 4)$ ist), während $a$, $c$, $d$ unverändert bleiben, so wächst $e$ um 2. Wir wollen die Gruppe weiter verfolgen:
\[
\begin{gathered}
n_{J+266}\equiv n_{J+42}\equiv n_{J+14}\equiv n_J+14+3\equiv n_J+3\; (\mathrm{mod}\ 7)\\
e_{J+266}\equiv e_J-3\equiv e_J+4\; (\mathrm{mod}\ 7)
\end{gathered}
\]
\noindent und
\[
\begin{gathered}
n_{J+399}\equiv n_{J+63}\equiv n_{J+7}\equiv n_J+7+1=n_J+1\; (\mathrm{mod}\ 7)\\
e_{J+399}\equiv e_J-1\equiv e_J+6\; (\mathrm{mod}\ 7).
\end{gathered}
\]
\noindent Ergebnis: $e$ ist eine lineare Funktion von $b$, und zwar ist
\[
e\equiv2b+\rho\; (\mathrm{mod}\ 7),
\]
\noindent wobei $\rho$ von $b$ unabhängig ist. Für $J+532$ treten bekanntlich die gleichen
Werte ein, wie für $J$, also ist auch $e_{J+532}=e_J$. $b$ ist um 3 zurückgegangen,
$e$ um 6, ganz im Sinne des eben aufgestellten Ergebnisses.

Endlich wählen wir die Gruppe:
\begin{center}
$\left.\begin{array}{l}J\\J+76\\J+152\\J+228\\J+304\\J+380\\J+456\end{array}\right\}$
\begin{minipage}{.81\linewidth}Auch hier hat $d$ für alle Jahreszahlen der Gruppe den gleichen
Wert. Der Siebenerrest $c$ nimmt jedesmal um 6 zu, oder, was
dasselbe heisst, um 1 ab. Nun ist
\[
n_{J+76}\equiv n_{J+20}\equiv n_J+20+5\equiv n_J+4\; (\mathrm{mod}\ 7),
\]
\noindent also
\[
e_{J+76}\equiv e_J-4\; (\mathrm{mod}\ 7).
\]
\end{minipage}
\end{center}
\clearpage\endgroup
